% Generated draft: requires mathematical and source review.
\documentclass[11pt]{article}
\usepackage[margin=25mm]{geometry}
\usepackage{fontspec}
\setmainfont{Latin Modern Roman}
\usepackage{amsmath,amssymb,mathtools,unicode-math}
\setmathfont{Latin Modern Math}
\usepackage{enumitem,needspace,xurl,hyperref}
\setlength{\parindent}{0pt}
\setlength{\parskip}{5pt}
\setlength{\emergencystretch}{3em}
\newcommand{\N}{\mathbb{N}}
\newcommand{\Z}{\mathbb{Z}}
\newcommand{\Q}{\mathbb{Q}}
\newcommand{\R}{\mathbb{R}}
\newcommand{\C}{\mathbb{C}}
\newcommand{\F}{\mathbb{F}}
\newcommand{\A}{\mathbb{A}}
\newcommand{\PP}{\mathbb{P}}
\newcommand{\E}{\mathbb{E}}
\newcommand{\eps}{\varepsilon}
\newcommand{\dd}{\,\mathrm{d}}
\begin{document}
\section*{UnsolvedMath ID: 701708}
\textbf{Source problem code:} catalogue entry
\subsection{Definitions and mathematical statement}
For a group \(G\), write
\[
I(G)=\{x\in G:x^2=1_G,\ x\ne 1_G\}
\]
for its set of involutions. For each positive integer \(m\), the source defines \(G_m\) as the ``largest'' group generated by \(m\) involutions and satisfying
\[
\forall x,y\in I(G_m),\qquad (xy)^4=1_{G_m}.
\]
Here ``largest'' is interpreted in the universal sense: \(G_m\) has distinguished involutory generators \(s_1,\ldots,s_m\), and for every group \(H=\langle t_1,\ldots,t_m\rangle\) with \(t_i\in I(H)\) and
\[
\forall u,v\in I(H),\qquad (uv)^4=1_H,
\]
there is a surjective homomorphism \(G_m\to H\) sending \(s_i\) to \(t_i\) for every \(i\). This makes explicit an interpretation of ``largest'' that the supplied record does not itself define. The condition concerns all involutions in the group, including those other than the distinguished generators.

\textbf{Problem.} Determine \(\lvert G_4\rvert\). No finiteness assumption is imposed here.

The source reports the known value \(\lvert G_3\rvert=2^{11}\) and refers additionally to Kourovka Notebook Problem 18.58; that referenced problem's contents are not supplied.
\subsection{Short English statement}
Determine the order of the universal group generated by four involutions in which the product of any two involutions has fourth power equal to the identity. The source reports \(\lvert G_3\rvert=2^{11}\); its term ``largest'' is interpreted universally and requires review.
\subsection{Sources}
\begin{itemize}
\item Supplied local UnsolvedMath registry record, permanent ID \(701708\), exact source code \texttt{catalogue entry}, titled \emph{Kourovka Notebook Problem 20.50}. The registry labels the problem open; this is supplied metadata, not an independently verified status assessment.
\item \emph{The Kourovka Notebook (Unsolved Problems in Group Theory)}, Problem 20.50, proposed by D. V. Lytkina in 2022. The supplied background identifies the 20th issue (2022) and the 21st issue edition (2026), \texttt{arXiv:1401.0300v48}. The supplied collection description names E. I. Khukhro and V. D. Mazurov as editors.
\item The supplied statement reports \(\lvert G_3\rvert=2^{11}\) and includes the cross-reference ``Also cf. 18.58.'' No separate reference supporting that numerical assertion or text of Problem 18.58 is supplied.
\end{itemize}
\end{document}
